% Propositions and proofs cited from Sections~\ref{sec:background}, \ref{sec:family}
% and \ref{sec:confound}. Input by appendix.tex under the section heading that
% carries \label{app:theory}, so this file starts at subsection level. Uses the
% `proposition' environment and amsthm's `proof', both set up in main.tex.

\subsection{Setting}
\label{app:theory-setting}

Fix a timestep $t$ and write $\sigma = \sigma_t$. We treat one channel; with
several channels every quantity below is averaged over channels, which leaves
the statements unchanged. The residual is $r \in \mathbb{R}^N$. A linear map
$W : \mathbb{R}^N \to \mathbb{R}^{N'}$ has its coefficients partitioned into $J$
bands $b \in \mathcal{B}$, with $P_b$ the restriction to the $n_b$ coefficients
of band $b$ and $\sum_b n_b = N'$. A complex unitary transform such as the DFT
is covered by the same arguments with the adjoint $W^{*}$ in place of
$W^\top$. For non-negative weights $\omega_b(\sigma)$
the per-sample loss of Eq.~\eqref{eq:family} and the pixel mean square are
\begin{equation}
  \ell_\sigma(r) = \sum_{b} \frac{\omega_b(\sigma)}{n_b} \lVert P_b W r \rVert_2^2
                 = r^\top Q_\sigma r,
  \qquad
  Q_\sigma = W^\top D_\sigma W,
  \qquad
  m(r) = \frac{\lVert r \rVert_2^2}{N},
  \label{eq:app-loss}
\end{equation}
where $D_\sigma$ is diagonal with entry $\omega_b(\sigma)/n_b$ on the
coefficients of band $b$. For $r \neq 0$ let
$\rho_b(r) = \operatorname{ms}(P_b W r)/m(r)$ be the mean square of band $b$
relative to the pixel mean square. If $W^\top W = I$, Eq.~\eqref{eq:parseval}
gives $\sum_b (n_b/N)\,\rho_b(r) = 1$; for orthogonal $W$ ($N' = N$),
$\rho_b \equiv 1$ when the residual's
energy is spread over the coefficients in proportion to band size (as for a
white residual in expectation).

\subsection{Band weighting leaves the minimiser unchanged}
\label{app:minimiser}

\begin{proposition}[Population minimiser]
\label{prop:minimiser}
Let $W$ be injective, and let $\omega_b(\sigma_t) > 0$ for every band and
timestep, with the weights depending on $(x_0, \epsilon, t)$ only through $t$.
For a measurable predictor $f$ let
$\mathcal{R}(f) = \mathbb{E}_{x_0,\epsilon,t}\big[\ell_{\sigma_t}\big(f(x_t,t) - \epsilon\big)\big]$.
Then $\mathcal{R}(f) \ge \mathcal{R}(f^\star)$ with
$f^\star(x_t, t) = \mathbb{E}[\epsilon \mid x_t, t]$, and equality holds only
if $f = f^\star$ almost everywhere. In particular, every such weighting,
including the uniform one that gives $\mathcal{L}_{\mathrm{simple}}$, has the
same minimiser.
\end{proposition}

\begin{proof}
For $r \neq 0$,
$r^\top Q_\sigma r = \sum_b (\omega_b/n_b)\lVert P_b W r\rVert_2^2
\ge (\min_b \omega_b/n_b)\,\lVert W r \rVert_2^2 > 0$, because $W r \neq 0$;
so $Q_\sigma$ is positive definite. Fix $(x_t, t)$ and write
$\mu = f^\star(x_t, t)$ and $y = f(x_t, t)$. Since $Q_{\sigma_t}$, $y$ and
$\mu$ are functions of $(x_t, t)$ and
$\mathbb{E}[\epsilon - \mu \mid x_t, t] = 0$, the cross term vanishes and
\begin{equation*}
  \mathbb{E}\big[(y-\epsilon)^\top Q_{\sigma_t} (y-\epsilon) \,\big|\, x_t, t\big]
  = (y-\mu)^\top Q_{\sigma_t} (y-\mu)
  + \mathbb{E}\big[(\epsilon-\mu)^\top Q_{\sigma_t} (\epsilon-\mu) \,\big|\, x_t, t\big].
\end{equation*}
The second term does not depend on $f$ and is finite because
$\mathbb{E}\lVert\epsilon\rVert^2 < \infty$; the first is non-negative and
vanishes only at $y = \mu$. Taking the expectation over $(x_t, t)$ completes
the proof.
\end{proof}

Three remarks delimit the result. (i) The argument is unchanged if $\ell$ is
multiplied by a positive timestep weight, so it covers the normalised
weightings of Section~\ref{sec:confound} and any product of a timestep
weighting with a band weighting. (ii) What the weighting does change is the
attained minimum (the second term above) and, when the model class does not
contain $f^\star$, which approximation is optimal: by the decomposition above,
the best predictor in a restricted class is the one closest to $f^\star$ in
the norm $\big(\mathbb{E}[(f - f^\star)^\top Q_{\sigma_t} (f - f^\star)]\big)^{1/2}$,
which depends on $\omega$. (iii) The proposition does not cover
weights that depend on $\epsilon$. If the weights are computed from the target,
as in the spatially adaptive arm of Section~\ref{sec:extensions}, whose map is
built from the target's coefficients, the loss is a quadratic form
$Q(\epsilon)$ and its minimiser is the weighted conditional mean
$\mathbb{E}[Q(\epsilon) \mid x_t, t]^{-1}\,\mathbb{E}[Q(\epsilon)\,\epsilon \mid x_t, t]$.
If the weights are computed from the residual and detached, as in FFL
(Eq.~\eqref{eq:ffl}), training does not descend a fixed objective; its
stationary points satisfy
$\mathbb{E}\big[F^{*}\operatorname{diag}(w(R))\,R \,\big|\, x_t, t\big] = 0$
with $R = F(f - \epsilon)$. Because $w$ is even in $R$, the conditional mean
satisfies this condition when the conditional law of $\epsilon$ is symmetric
about its mean, but not in general.

\subsection{The timestep component of a band weighting}
\label{app:confound}

\begin{proposition}[Timestep and shape decomposition]
\label{prop:confound}
Let $W$ be orthogonal ($N' = N$ and $W^\top W = I$) and let the weights
$\omega_b(\sigma) \ge 0$ depend on $\sigma$ only. Write $S(\sigma) = \sum_b \omega_b(\sigma)$ and
$\delta_b(\sigma) = \omega_b(\sigma) - (n_b/N)\,S(\sigma)$, so that
$\sum_b \delta_b = 0$. Then:
\begin{enumerate}
  \item[(a)] For every residual $r$,
  \begin{equation}
    \ell_\sigma(r)
    = S(\sigma)\, m(r)
    + \sum_b \delta_b(\sigma)\,\big[\operatorname{ms}(P_b W r) - m(r)\big]
    = m(r) \sum_b \omega_b(\sigma)\,\rho_b(r) .
    \label{eq:shape-decomp}
  \end{equation}
  The second term vanishes identically when $\omega_b \propto n_b$, and for a
  given residual whenever $\rho_b(r) = 1$ for every band with
  $\delta_b \neq 0$.
  \item[(b)] Let
  $\bar\rho_b(t) = \mathbb{E}[\operatorname{ms}(P_b W r) \mid t]/\mathbb{E}[m(r) \mid t]$.
  Then $\mathbb{E}[\ell_{\sigma_t}(r) \mid t] = \Lambda(t)\,\mathbb{E}[m(r) \mid t]$ with
  \begin{equation}
    \Lambda(t) = \sum_b \omega_b(\sigma_t)\,\bar\rho_b(t)
    = S(\sigma_t) + \sum_b \delta_b(\sigma_t)\,\big(\bar\rho_b(t) - 1\big)
    = \sum_b \frac{N}{n_b}\,\omega_b(\sigma_t)\, e_b(t),
    \label{eq:eff-weight}
  \end{equation}
  where $e_b(t) = (n_b/N)\,\bar\rho_b(t)$ are the residual's band energy
  fractions at timestep $t$ ($\sum_b e_b = 1$). The effective timestep weight
  $\Lambda(t)$ is thus a convex combination of the pixel-equivalent weights
  $(N/n_b)\,\omega_b(\sigma_t)$, and lies between their minimum and maximum.
  \item[(c)] For AWWL (Eq.~\eqref{eq:audited}; $J = 4$, $n_b = N/4$),
  $S(\sigma) = (1-\alpha) + (2\alpha-1)\lambda_p(\sigma)$ and
  \begin{align}
    \ell_\sigma(r) &= S(\sigma)\, m(r)
      + \Big(w_{\mathrm{LL}}(\sigma) - \frac{w_{\det}(\sigma)}{3}\Big)
        \big[\operatorname{ms}(r_{\mathrm{LL}}) - m(r)\big],
    \label{eq:awwl-decomp}\\
    \Lambda(t) &= \frac{4}{3}\,w_{\det}(\sigma_t)
      + \Big(w_{\mathrm{LL}}(\sigma_t) - \frac{w_{\det}(\sigma_t)}{3}\Big)\,\bar\rho_{\mathrm{LL}}(t),
    \label{eq:awwl-eff}
  \end{align}
  with $\bar\rho_{\mathrm{LL}}(t) \in [0, 4]$. The residual-independent weight
  $S$ is constant in $\sigma$ if and only if $\alpha = 1/2$ (for $p > 0$). For
  $\alpha \neq 1/2$ and any function $\phi$ of the timestep, the correlation over
  timesteps satisfies
  $\operatorname{corr}_t\big(S(\sigma_t), \phi(t)\big)
  = \operatorname{sign}(2\alpha - 1)\,\operatorname{corr}_t\big(\lambda_p(\sigma_t), \phi(t)\big)$,
  independently of $\lvert\alpha - 1/2\rvert$.
  \item[(d)] Rescaling the AWWL weights to a constant total $c$, that is,
  replacing $\omega_b$ by $\tilde\omega_b = c\,\omega_b/S$, removes the
  residual-independent timestep weight exactly and leaves
  $\tilde\Lambda(t) = c + c\,\big(w_{\mathrm{LL}} - w_{\det}/3\big)\big(\bar\rho_{\mathrm{LL}}(t) - 1\big)/S$,
  evaluated at $\sigma_t$. It equals $c$ at every timestep if and only if, at
  every timestep, either the approximation band carries a quarter of the
  residual's energy ($\bar\rho_{\mathrm{LL}}(t) = 1$) or the four per-band
  weights are equal ($w_{\mathrm{LL}} = w_{\det}/3$, which holds only at
  $\lambda_p(\sigma) = (1-\alpha)/(1+2\alpha)$).
\end{enumerate}
\end{proposition}

\begin{proof}
(a) Since $W^\top W = I$, Eq.~\eqref{eq:parseval} gives
$\sum_b n_b \operatorname{ms}(P_b W r) = \lVert r \rVert_2^2 = N m(r)$. Hence
\[
  \textstyle\sum_b \delta_b \operatorname{ms}(P_b W r)
  = \ell_\sigma(r) - (S/N)\sum_b n_b \operatorname{ms}(P_b W r)
  = \ell_\sigma(r) - S\,m(r),
\]
and $\sum_b \delta_b\, m(r) = 0$, which gives the first equality; the second is
the definition of $\rho_b$. If
$\omega_b \propto n_b$ then $\delta_b = 0$ for all $b$.
(b) Take conditional expectations in Eq.~\eqref{eq:shape-decomp} at timestep
$t$ and divide by $\mathbb{E}[m \mid t]$. The last form uses
$\omega_b \bar\rho_b = (N/n_b)\,\omega_b\, e_b$; $e_b \ge 0$ and
$\sum_b e_b = 1$ by Eq.~\eqref{eq:parseval}.
(c) With four bands of size $N/4$, $\sum_b \operatorname{ms}(r_b) = 4\,m(r)$, so
$\sum_{b \in \{\mathrm{LH},\mathrm{HL},\mathrm{HH}\}} \big(\operatorname{ms}(r_b) - m\big)
= -\big(\operatorname{ms}(r_{\mathrm{LL}}) - m\big)$. The three detail bands
share $\delta_{\det} = w_{\det}/3 - S/4$, and
$\delta_{\mathrm{LL}} - \delta_{\det} = w_{\mathrm{LL}} - w_{\det}/3$, so the
second term of Eq.~\eqref{eq:shape-decomp} equals
$(\delta_{\mathrm{LL}} - \delta_{\det})\big(\operatorname{ms}(r_{\mathrm{LL}}) - m\big)$.
Substituting $S = w_{\mathrm{LL}} + w_{\det}$ gives the expression for
$\Lambda$. The bound $\bar\rho_{\mathrm{LL}} \le 4$ follows from
$\lVert r_{\mathrm{LL}} \rVert^2 \le \lVert r \rVert^2$. $S$ is affine in
$\lambda_p$ with slope $2\alpha - 1$ by Eq.~\eqref{eq:total}, and $\lambda_p$
is strictly increasing on $(0,1)$ for $p > 0$; a Pearson correlation is
invariant under increasing affine maps of either argument and changes sign
under decreasing ones.
(d) Replace $\omega_b$ by $c\,\omega_b/S$ in (c).
\end{proof}

In Eq.~\eqref{eq:shape-decomp} the first term is the pixel MSE weighted by
$S(\sigma_t)$, a timestep weight that does not depend on the residual. The
second applies the zero-sum weights $\delta_b$ to the departures of the band
mean squares from the pixel mean square; it carries the frequency shape of the
weighting. By Eq.~\eqref{eq:family-x0}, the weight that the objective places on
the clean-image error at timestep $t$ is $\mathrm{SNR}_t\,\Lambda(t)$. Parts
(b) and (d) locate the part of the timestep weighting that no rescaling of the
weights by a function of $\sigma$ fixed in advance removes: the interaction of the zero-sum part $\delta_b$ with the
residual's band energy fractions, which change with $\sigma_t$
(Proposition~\ref{prop:residual}).

\subsection{What the optimiser receives}
\label{app:scale}

\begin{proposition}[Loss scale and gradient scale]
\label{prop:scale}
\begin{enumerate}
  \item[(a)] Consider the AdamW update
  \begin{equation}
    \theta_{k+1} = (1 - \eta_k\zeta)\,\theta_k
      - \eta_k\,\frac{\hat m_k}{\sqrt{\hat v_k} + \varepsilon_{\mathrm{A}}},
    \label{eq:adamw}
  \end{equation}
  where $\eta_k$ is the learning rate, $\zeta$ the decoupled weight decay,
  $\hat m_k$ and $\hat v_k$ the bias-corrected exponential moving averages of
  the gradients $g_k$ and of $g_k \odot g_k$, and the square root and division
  act elementwise. If every gradient is multiplied by a constant $c > 0$, the
  iterates are those of the unscaled run with $\varepsilon_{\mathrm{A}}$
  replaced by $\varepsilon_{\mathrm{A}}/c$. With gradient-norm clipping at a
  threshold $\tau$ applied before the update, the same holds if clipping is
  inactive on every step of both runs, and the two runs coincide exactly if it
  is active on every step of both.
  \item[(b)] For one sample at noise level $\sigma$,
  $\nabla_\theta \ell_\sigma = \sum_b \omega_b(\sigma)\, q_b$ with
  $q_b = \nabla_\theta \operatorname{ms}(P_b W r)$. If $\lVert q_b \rVert = G$ for
  every band, then $\lVert \nabla_\theta \ell_\sigma \rVert \le G\,S(\sigma)$, with
  equality when the $q_b$ are positively parallel. If in addition
  $\langle q_b, q_c \rangle \ge 0$ for all $b, c$, then
  $\lVert \nabla_\theta \ell_\sigma \rVert \ge G\,\big(\sum_b \omega_b(\sigma)^2\big)^{1/2}$,
  with equality when the $q_b$ are mutually orthogonal.
  \item[(c)] Let $W$ be orthogonal ($N' = N$, $W^\top W = W W^\top = I$). The
  output-space band gradients
  $h_b = \nabla_{\hat\epsilon} \operatorname{ms}(P_b W r) = (2/n_b)\, W^\top P_b^\top P_b W r$
  are mutually orthogonal, so
  \begin{equation}
    \lVert \nabla_{\hat\epsilon}\, \ell_\sigma \rVert_2^2
    = \sum_b \Big(\frac{2\,\omega_b(\sigma)}{n_b}\Big)^2 \lVert P_b W r \rVert_2^2 ,
    \label{eq:gradnorm}
  \end{equation}
  and for $J$ bands of equal size with $\rho_b(r) = 1$,
  $\lVert \nabla_{\hat\epsilon}\, \ell_\sigma \rVert_2 = (2\sqrt{J}/N)\,\lVert r \rVert_2\,\big(\sum_b \omega_b^2\big)^{1/2}$.
  For every $W$, the component of $\nabla_{\hat\epsilon}\, \ell_\sigma$ along
  the pixel-MSE gradient $\nabla_{\hat\epsilon}\, m$ has coefficient
  \begin{equation}
    \frac{\langle \nabla_{\hat\epsilon}\, \ell_\sigma, \nabla_{\hat\epsilon}\, m \rangle}
         {\lVert \nabla_{\hat\epsilon}\, m \rVert_2^2}
    = \frac{\ell_\sigma(r)}{m(r)} = \sum_b \omega_b(\sigma)\,\rho_b(r),
    \label{eq:mse-component}
  \end{equation}
  which equals $S(\sigma)$ when $\rho_b(r) = 1$ for every band.
  \item[(d)] For $\omega \in [0,\infty)^J$ with $S = \sum_b \omega_b$,
  \begin{equation}
    S^2/J \;\le\; \textstyle\sum_b \omega_b^2 \;\le\; S^2 ,
    \label{eq:cs-bounds}
  \end{equation}
  with equality on the left if and only if the weights are equal, and on the
  right if and only if at most one weight is non-zero. Holding $S$ fixed
  therefore leaves $\sum_b \omega_b^2$ free within a factor $J$, and holding
  $\sum_b \omega_b^2$ fixed leaves $S$ free within a factor $\sqrt{J}$.
\end{enumerate}
\end{proposition}

\begin{proof}
(a) $\hat m_k$ is linear in $g_1, \dots, g_k$ and $\hat v_k$ in
$g_1 \odot g_1, \dots, g_k \odot g_k$, so they scale by $c$ and $c^2$, and
$c\,\hat m_k/(c\sqrt{\hat v_k} + \varepsilon_{\mathrm{A}})
= \hat m_k/(\sqrt{\hat v_k} + \varepsilon_{\mathrm{A}}/c)$. The weight-decay term
does not involve the gradient, and the claim follows by induction on $k$.
Clipping multiplies $g$ by $\min(1, \tau/\lVert g \rVert)$; if it is inactive on
every step it changes nothing, and if it is active on every step both runs
receive $\tau\, g/\lVert g \rVert$, which does not depend on $c$.
(b) The upper bound is the triangle inequality. For the lower bound,
$\lVert \sum_b \omega_b q_b \rVert^2
= G^2 \sum_b \omega_b^2 + \sum_{b \neq c} \omega_b \omega_c \langle q_b, q_c \rangle
\ge G^2 \sum_b \omega_b^2$, since $\omega \ge 0$.
(c) $\nabla_{\hat\epsilon}\, \ell_\sigma = 2 Q_\sigma r = \sum_b \omega_b h_b$. For
$b \neq c$,
$\langle h_b, h_c \rangle \propto r^\top W^\top P_b^\top P_b\, W W^\top P_c^\top P_c W r = 0$,
because $W W^\top = I$ and $P_b P_c^\top = 0$; and
$\lVert h_b \rVert^2 = (4/n_b^2)\lVert P_b W r \rVert^2$, because $W^\top$ and
$P_b^\top$ preserve norms. With $n_b = N/J$ and
$\lVert P_b W r \rVert^2 = \lVert r \rVert^2/J$, Eq.~\eqref{eq:gradnorm} gives the
stated norm. Finally $\nabla_{\hat\epsilon}\, m = 2r/N$, and
$\langle 2 Q_\sigma r, 2r/N \rangle / \lVert 2r/N \rVert^2
= N\, r^\top Q_\sigma r / \lVert r \rVert^2 = \ell_\sigma(r)/m(r)$.
(d) The left inequality is the Cauchy--Schwarz inequality
$(\sum_b 1 \cdot \omega_b)^2 \le J \sum_b \omega_b^2$. The right one follows
from $S^2 = \sum_b \omega_b^2 + \sum_{b \neq c} \omega_b \omega_c$ with every cross
product non-negative, and it is an equality exactly when every cross product
vanishes.
\end{proof}

Part (c) is the orthogonal case of part (b): in output space the band
gradients are orthogonal, and with equal band energies the gradient norm scales
with $(\sum_b \omega_b^2)^{1/2}$. The parameter gradient is their image under
the transposed network Jacobian, which can align them; in the extreme case the
images are parallel and the norm scales with $S$. When the images have equal
norms and non-negative inner products (the assumptions of part (b)), the
per-timestep gradient contribution therefore lies between the two cases, and
the two normalisations,
$S$ held fixed and $\sum_b \omega_b^2$ held fixed, are the natural bounding
controls: the first fixes the contribution in the parallel case (and, in output
space, the component along the pixel-MSE gradient), the second fixes it in the
orthogonal case. By part (d) no rescaling fixes both unless the weights are
uniform, and both statements assume $\rho_b = 1$, the case in which the
residual-dependent term of Proposition~\ref{prop:confound} vanishes. By part
(a), neither control depends on the value of the constant it holds fixed,
except through $\varepsilon_{\mathrm{A}}$ and clipping; only the variation of
the weights across timesteps and bands matters.

For AWWL at $\alpha = 0.2$, $p = 1$ on the CIFAR-10 schedule, holding $S$
fixed leaves $\sum_b \omega_b^2$ varying by a factor of $4.00 = J$ over the
training timesteps. Both bounds of Eq.~\eqref{eq:cs-bounds} are approached:
the left one at $\lambda_p(\sigma) = (1-\alpha)/(1+2\alpha)$, that is
$\sigma = 4/7$, where the four per-band weights are equal, and the right one as
$\sigma \to 1$, where all weight moves to the approximation band.

\subsection{The residual spectrum under a Gaussian model}
\label{app:residual}

\begin{proposition}[Residual of the optimal predictor]
\label{prop:residual}
Let $x_0 \sim \mathcal{N}(0, C)$ with $C = U \operatorname{diag}(P)\, U^\top$ for
an orthonormal basis $U$. For a stationary (circulant) covariance, $U$ can be
taken to be the real Fourier basis and $P$ is the power spectrum. Let
$\hat\epsilon = \mathbb{E}[\epsilon \mid x_t]$ be the optimal predictor at
timestep $t$ (Proposition~\ref{prop:minimiser}). Then the coordinates of
$U^\top(\hat\epsilon - \epsilon)$ are independent, with variances
\begin{equation}
  v_t(u) = \frac{\bar\alpha_t P(u)}{\bar\alpha_t P(u) + \sigma_t^2} .
  \label{eq:gaussian-residual}
\end{equation}
For two coordinates with $P(u) > P(u') > 0$, the ratio $v_t(u)/v_t(u')$
increases strictly from $1$ to $P(u)/P(u')$ as $\mathrm{SNR}_t$ decreases from
$\infty$ to $0$.
\end{proposition}

\begin{proof}
Write $z = U^\top x_0$, $\tilde\epsilon = U^\top \epsilon$ and
$y = U^\top x_t = \sqrt{\bar\alpha_t}\, z + \sigma_t \tilde\epsilon$. Because $U$ is
orthonormal, $\tilde\epsilon \sim \mathcal{N}(0, I)$, and $z$ and $\tilde\epsilon$ are independent with
independent coordinates, so the problem decouples across coordinates. For one
coordinate, $(\tilde\epsilon_u, y_u)$ is jointly Gaussian with
$\operatorname{Var}(y_u) = \bar\alpha_t P(u) + \sigma_t^2$ and
$\operatorname{Cov}(\tilde\epsilon_u, y_u) = \sigma_t$, so
$\mathbb{E}[\tilde\epsilon_u \mid y] = \sigma_t\, y_u/(\bar\alpha_t P(u) + \sigma_t^2)$ and
the residual variance is
$1 - \sigma_t^2/(\bar\alpha_t P(u) + \sigma_t^2) = v_t(u)$; each residual
coordinate depends only on $(z_u, \tilde\epsilon_u)$, hence independence. With
$s = 1/\mathrm{SNR}_t = \sigma_t^2/\bar\alpha_t$,
$v_t(u)/v_t(u') = \big(P(u)/P(u')\big)\,\big(P(u') + s\big)/\big(P(u) + s\big)$,
whose derivative in $s$ has the sign of $P(u) - P(u') > 0$; it equals $1$ at
$s = 0$ and tends to $P(u)/P(u')$ as $s \to \infty$.
\end{proof}

For natural images $P$ decreases with frequency, so the optimal predictor's
residual is nearly white at low noise ($v_t \to 1$) and shifts its energy
towards low frequencies as the noise level grows. The normalised residual
spectra $v_t/\sum_u v_t(u)$ are ordered in the likelihood-ratio sense in $P$
(for $\sigma_{t'} > \sigma_t$ the ratio $v_{t'}(u)/v_t(u)$ increases with
$P(u)$), so the share of residual energy in any band whose squared frequency
response is an increasing function of $P$ grows with $\sigma_t$. The Haar
approximation band is not exactly of this form, because its response is not a
function of radial frequency alone, but its response is concentrated at the
low frequencies where a natural-image spectrum is largest; evaluating
$\bar\rho_{\mathrm{LL}}(t)$ directly for $P(u) \propto \lvert u \rvert^{-2}$ on
the CIFAR-10 schedule gives a monotone increase from $1$ at the lowest noise
level. In Proposition~\ref{prop:confound} this makes $\Lambda(t)$ depend on the
timestep even when $S$ is held fixed. How far a trained model's residual
follows this pattern is an empirical question.

\subsection{Gradient-norm bounds}
\label{app:gradnorm}

The two controls of Section~\ref{sec:confound-controls} hold different
quantities fixed because the loss scale and the gradient scale of a band
weighting differ. For one sample,
$\nabla_\theta \ell_\sigma = \sum_b \omega_b(\sigma)\, q_b$, with $q_b$ the
parameter gradient of the mean square of band $b$. If the $q_b$ have a common
norm $G$ and non-negative inner products, Proposition~\ref{prop:scale}(b)
gives
\begin{equation}
  G \Big(\sum_b \omega_b(\sigma)^2\Big)^{1/2}
  \;\le\; \lVert \nabla_\theta \ell_\sigma \rVert
  \;\le\; G\, S(\sigma).
  \label{eq:gradscale}
\end{equation}
The upper bound holds with equality when the band gradients are parallel, the
lower when they are orthogonal, as the output-space band gradients are
(Proposition~\ref{prop:scale}(c)). Holding $S$ fixed therefore fixes the
per-timestep gradient weight in the parallel case and, for residuals with
equal band energies, the output-space component along the pixel-MSE gradient
(Eq.~\eqref{eq:mse-component}); holding $\sum_b \omega_b^2$ fixed fixes it in
the orthogonal case. By Eq.~\eqref{eq:cs-bounds} with $J = 4$ bands, fixing
either quantity leaves the other free within a factor of $\sqrt{J} = 2$ in
norm, and both limits are approached on the CIFAR-10 schedule
(Fig.~\ref{fig:gradient-scale}; the points at which they are reached are given
at the end of Appendix~\ref{app:scale}). Neither control is therefore the
stronger: each fixes one end of Eq.~\eqref{eq:gradscale}. The first control
multiplies the AWWL loss by $c/S(\sigma)$ and the second by
$\big(\levelSqFixed / \sum_b \omega_b(\sigma)^2\big)^{1/2}$. Both statements
about the gradient weight assume $\rho_b(r) = 1$; for other residuals the
residual-dependent term of Proposition~\ref{prop:confound} remains, and
neither control fixes the effective timestep weight $\Lambda(t)$.

\guardedfigure{1.0}{gradient_scale.png}{%
  What each control holds fixed, for the released form of AWWL at
  $\alpha = 0.2$, $p = 1$, with per-band weights
  $\omega_b = (w_{\mathrm{LL}}, w_{\det}/3, w_{\det}/3, w_{\det}/3)$ as in
  the panel titles. (a) The total
  $S = \sum_b \omega_b$, which sets the loss scale at each timestep and the
  gradient weight when the band gradients are parallel. (b)
  $\sum_b \omega_b^2$, which sets the squared gradient weight when the band
  gradients are orthogonal with equal norms (Eq.~\eqref{eq:gradscale}). Each
  curve is divided by AWWL's mean of the same quantity over the training
  timesteps (\wTotalMean{} in (a), \wSqMean{} in (b)). Solid: AWWL; dotted:
  total fixed at \levelSumFixed; dash-dotted: sum of squares fixed at
  \levelSqFixed. Fixing either quantity leaves the other varying, by a factor
  of $4$ in (b) (of $2$ in its square root) and of $2$ in (a). Horizontal
  axes: training timestep $t$, sampled uniformly; top axes:
  $\sigma_t$.}{fig:gradient-scale}
